\DocumentMetadata{
  pdfversion=2.0,
  pdfstandard=ua-2,
  lang=en-US,
  tagging=on,
  tagging-setup={math/setup=mathml-SE,tabsorder=structure}
}
\documentclass[11pt,letterpaper]{article}
\usepackage[margin=0.68in,includefoot]{geometry}
\usepackage{newcomputermodern}
\usepackage{amsmath,amscd,mathtools}
\usepackage{tikz,adjustbox}
\providecommand{\square}{\mdlgwhtsquare}
\usepackage[hidelinks]{hyperref}
\hypersetup{
  pdftitle={Algebra PhD Exam, August 2023},
  pdfauthor={University of Florida Department of Mathematics},
  pdfsubject={Graduate mathematics examination},
  pdfdisplaydoctitle=true,
  bookmarksopen=true
}
\setcounter{secnumdepth}{0}
\setlength{\parindent}{0pt}
\setlength{\parskip}{0.28em}
\setlength{\emergencystretch}{3em}
\setlength{\leftmargini}{2.15em}
\setlength{\leftmarginii}{2.65em}
\setlength{\leftmarginiii}{3em}
\renewcommand{\labelenumi}{\textbf{\theenumi.}}
\pagestyle{empty}
\raggedbottom
\allowdisplaybreaks
\newenvironment{examproblems}[1]{%
  \begin{enumerate}%
  \setcounter{enumi}{#1}%
  \setlength{\topsep}{0.35em}%
  \setlength{\partopsep}{0pt}%
  \setlength{\itemsep}{0.48em}%
  \setlength{\parsep}{0.18em}%
}{\end{enumerate}}
\newenvironment{examsubparts}{%
  \begin{enumerate}%
  \setlength{\topsep}{0.18em}%
  \setlength{\partopsep}{0pt}%
  \setlength{\itemsep}{0.16em}%
  \setlength{\parsep}{0.1em}%
}{\end{enumerate}}
\newenvironment{examinstructions}{%
  \par\begingroup\itshape
}{%
  \par\endgroup\vspace{0.18em}
}
\makeatletter
\renewcommand\section{\@startsection{section}{1}{0pt}%
  {-1.25ex plus -.25ex minus -.1ex}%
  {0.55ex plus .1ex}%
  {\normalfont\large\bfseries}}
\renewcommand{\@maketitle}{%
  \newpage\null\vskip -1em%
  {\footnotesize\bfseries
    \href{https://www.ufl.edu/}{University of Florida}, \href{https://math.ufl.edu/}{Department of Mathematics}\par}%
  \vskip -0.5em%
  \tagstructbegin{tag=Title}%
    {\LARGE\bfseries\href{https://gma.math.ufl.edu/exams/algebra-phd/}{Algebra PhD Exam}, August 2023\par}%
  \tagstructend%
  \vskip 0.25em%
  \tagmcbegin{artifact}%
    \hrule height 1pt%
  \tagmcend%
  \vskip 1em%
}
\makeatother
\title{Algebra PhD Exam, August 2023}
\author{}
\date{}
\begin{document}
\maketitle
\thispagestyle{empty}
\begin{examinstructions}
Answer seven problems. Write your answers clearly in complete English sentences. You may quote results (within reason) as long as you state them clearly.
\end{examinstructions}
\begin{examproblems}{0}
\item
Briefly explain your answers to the following, freely using standard results.
\begin{examsubparts}
\item[\textnormal{(a)}]
Find a polynomial \(f(x)\in\mathbb{F}_3[x]\) such that \(K=\mathbb{F}_3[x]/(f(x))\) is a field with 27 elements.
\item[\textnormal{(b)}]
Show that~\(K\) is Galois over \(\mathbb{F}_3\), and write down the automorphisms of~\(K\) fixing \(\mathbb{F}_3\).
\item[\textnormal{(c)}]
For which~\(n\) is there an extension \(L/K\) of fields such that~\(L\) has \(3^n\) elements?
\end{examsubparts}
\item
Let \(f(x)\in\mathbb{Q}[x]\) be an irreducible polynomial of degree 5 with exactly three real roots. State and prove a result about the Galois group of the splitting field of \(f(x)\) over~\(\mathbb{Q}\).
\item
Let \(K=\mathbb{Q}(\zeta_{101})\), where \(\zeta_{101}\) is a primitive 101th root of unity.
\begin{examsubparts}
\item[\textnormal{(a)}]
Show that there is a unique subfield~\(F\) of~\(K\) such that \([F:\mathbb{Q}]=2\).
\item[\textnormal{(b)}]
How many subfields of~\(K\) contain~\(F\)? Of these, which are Galois extensions of~\(\mathbb{Q}\)?
\end{examsubparts}
\item
Recall that the Abelianization of a group~\(G\) is \(G^{\mathrm{ab}}:=G/[G,G]\).
\begin{examsubparts}
\item[\textnormal{(a)}]
Formulate and prove a universal property for abelianization.
\item[\textnormal{(b)}]
Show that abelianization is left adjoint to the forgetful functor from the category of abelian groups to the category of groups.
\end{examsubparts}
\item
Let~\(K\) be a field and~\(V\) a~\(K\)-vector space. Recall that an alternating bilinear map \(f:V\times V\to W\) is a function which is~\(K\)-linear in each coordinate (i.e. \(K\)-multilinear) and such that \(f(v,v)=0\) for all \(v\in V\). Show that the functor \(\operatorname{Alt}_V^2:\operatorname{Vec}_K\to\operatorname{Sets}\) defined by 
\[
\operatorname{Alt}_V^2(W)=\{f:V\times V\to W:f\text{ is alternating bilinear}\}
\]
 is a representable functor. (The representing object, usually denoted \(\Lambda^2(V)\), is an exterior power of~\(V\).) Suggestion: A quotient of \(V\otimes_K V\).
\item
Give examples of the following, and briefly explain your answers.
\begin{examsubparts}
\item[\textnormal{(a)}]
a ring~\(R\) and a short exact sequence of~\(R\)-modules 
\[
0\to A\to B\to C\to 0
\]
 that is not split.
\item[\textnormal{(b)}]
a flat~\(\mathbb{Z}\)-module which is not free.
\item[\textnormal{(c)}]
a non-zero injective~\(\mathbb{Z}\)-module.
\end{examsubparts}
\item
Let~\(R\) be a ring and~\(M\) and~\(N\) be~\(R\)-modules. Let \(P_\bullet\) and \(P'_\bullet\) be projective resolutions of~\(M\) and~\(N\). Given a morphism of~\(R\)-modules \(f:M\to N\), show that there exists a morphism of resolutions \(f:P_\bullet\to P'_\bullet\) extending~\(f\). In other words, show there exists \(f_0,f_1,f_2,\ldots\) making the following diagram commute: 
{\tagpdfsetup{math/mathml/structelem=false,math/mathml/AF=false,math/alt/use=true}
\[
\begin{CD}
\cdots @>>> P_2 @>{d_2}>> P_1 @>{d_1}>> P_0 @>{\epsilon}>> M \\
@. @V{f_2}VV @V{f_1}VV @V{f_0}VV @V{f}VV \\
\cdots @>>> P_2' @>{d_2'}>> P_1' @>{d_1'}>> P_0' @>{\epsilon'}>> N
\end{CD}
\]
}
\item
What are the dimensions of \(\mathbb{Q}[x]/(x^5)\otimes_{\mathbb{Q}}\mathbb{Q}[x]/(x^6)\) and of \(\mathbb{Q}[x]/(x^5)\otimes_{\mathbb{Q}[x]}\mathbb{Q}[x]/(x^6)\) as~\(\mathbb{Q}\)-vector spaces?
\item
Let~\(F\) be a field. Prove that the power series ring 
\[
F[\![x]\!]=\left\{\sum_{n=0}^{\infty}a_nx^n:a_n\in F\right\}
\]
 is a Noetherian ring.
\item
Let \(B=\mathbb{Q}[x,y]/(xy)\) and \(A=\mathbb{Q}[t]\). Note that both have Krull dimension 1. A homomorphism \(f:A\to B\) of rings turns~\(B\) into an~\(A\)-module: for \(a\in A\) and \(b\in B\), \(a.b=f(a)b\).
\begin{examsubparts}
\item[\textnormal{(a)}]
Explain why the Noether normalization lemma implies that there is a map \(A\to B\) making~\(B\) into a finitely generated~\(A\)-module.
\item[\textnormal{(b)}]
Show that the map \(A\to B\) sending~\(t\) to~\(x\) does not make~\(B\) into a finitely generated~\(A\)-module.
\item[\textnormal{(c)}]
Find an explicit map \(A\to B\) that makes~\(B\) into a finitely generated~\(A\)-module and explain your answer.
\end{examsubparts}
\item
Let~\(K\) be a field and consider the function \(v:K(x)\to\mathbb{Z}\cup\{\infty\}\) given by 
\[
v(f/g)=\deg(g)-\deg(f),\qquad v(0)=\infty.
\]
\begin{examsubparts}
\item[\textnormal{(a)}]
Show that~\(v\) is a discrete valuation on \(K(x)\).
\item[\textnormal{(b)}]
Show that the valuation ring \(\mathcal{O}_v=\{R\in K(x):v(R)\geq 0\}\) is a local ring. (Do not appeal to general facts about DVR’s.)
\end{examsubparts}
\end{examproblems}
\end{document}
